\textbf{Phần 1. Trắc nghiệm lựa chọn}
\setcounterpageref{ex}{0}

\begin{ex} %[3D1V46]
Điền dấu thích hợp vào chỗ chấm: $4582 \dots 4590$.
\choice
{$>$}
{\True $<$}
{$=$}
{Không so sánh được}
\loigiai{
Hàng nghìn và hàng trăm bằng nhau, so sánh hàng chục: $8 < 9$ nên $4582 < 4590$.
\textbf{Đáp án của câu hỏi là B}
}
\end{ex}

\begin{ex} %[3D1V46]
Trong các số $7856, 7865, 8756, 8675$, số lớn nhất là
\choice
{$7865$}
{$8675$}
{$7856$}
{\True $8756$}
\loigiai{
So sánh hàng nghìn: số 8756 và 8675 lớn hơn. Tiếp tục so sánh hàng trăm: $7 > 6$ nên $8756$ lớn nhất.
\textbf{Đáp án của câu hỏi là D}
}
\end{ex}

\begin{ex} %[3D1V46]
Tìm chữ số thích hợp thay vào dấu $*$ để: $68*4 > 6884$. Chữ số đó là
\choice
{\True $9$}
{$8$}
{$7$}
{$6$}
\loigiai{
Hàng nghìn, hàng trăm và hàng đơn vị bằng nhau. Để $68*4 > 6884$ thì $* > 8$, suy ra $* = 9$.
\textbf{Đáp án của câu hỏi là A}
}
\end{ex}

\begin{ex} %[3D1V46]
Sắp xếp các số sau theo thứ tự từ bé đến lớn: $3050, 3500, 3005, 3505$.
\choice
{$3505, 3500, 3050, 3005$}
{$3050, 3005, 3500, 3505$}
{\True $3005, 3050, 3500, 3505$}
{$3005, 3500, 3050, 3505$}
\loigiai{
So sánh các số: $3005 < 3050 < 3500 < 3505$.
\textbf{Đáp án của câu hỏi là C}
}
\end{ex}

\begin{ex} %[3D1V46]
Có bao nhiêu số tự nhiên $x$ có bốn chữ số thỏa mãn: $2998 < x < 3003$?
\choice
{$3$ số}
{\True $4$ số}
{$5$ số}
{$6$ số}
\loigiai{
Các số tự nhiên $x$ là: $2999, 3000, 3001, 3002$. Có 4 số thỏa mãn.
\textbf{Đáp án của câu hỏi là B}
}
\end{ex}

\begin{ex} %[3D1V46]
Số bé nhất trong các số sau là số nào: $999, 1002, 1020, 2001$?
\choice
{$1002$}
{$1020$}
{$2001$}
{\True $999$}
\loigiai{
Số có ít chữ số hơn thì bé hơn. Số 999 có 3 chữ số nên bé nhất.
\textbf{Đáp án của câu hỏi là D}
}
\end{ex}

\begin{ex} %[3D1V46]
Tìm số tự nhiên tròn trăm $x$ thỏa mãn: $4500 < x < 4800$. Các giá trị của $x$ là
\choice
{\True $4600$ và $4700$}
{$4600$}
{$4700$}
{$4550$ và $4650$}
\loigiai{
Số tròn trăm nằm giữa 4500 và 4800 là 4600 và 4700.
\textbf{Đáp án của câu hỏi là A}
}
\end{ex}

\begin{ex} %[3D1V46]
So sánh hai biểu thức: $A = 3000 + 400 + 50 + 2$ và $B = 3450 + 5$. Khẳng định đúng là
\choice
{$A > B$}
{$A = B$}
{\True $A < B$}
{Không so sánh được}
\loigiai{
$A = 3452$; $B = 3455$. Vì $3452 < 3455$ nên $A < B$.
\textbf{Đáp án của câu hỏi là C}
}
\end{ex}

\begin{ex} %[3D1V46]
Bốn bạn An, Bình, Cúc, Dũng thi chạy. Điểm tích lũy lần lượt là: An được 4520 điểm, Bình được 4250 điểm, Cúc được 4502 điểm, Dũng được 4052 điểm. Bạn đạt điểm cao nhất là
\choice
{Bình}
{\True An}
{Cúc}
{Dũng}
\loigiai{
So sánh các số: $4520 > 4502 > 4250 > 4052$. Điểm cao nhất là An (4520 điểm).
\textbf{Đáp án của câu hỏi là B}
}
\end{ex}

\begin{ex} %[3D1V46]
Tìm chữ số $x$ biết rằng: $5x78 < 5178$. Chữ số $x$ là
\choice
{$1$}
{$2$}
{Không có chữ số nào}
{\True $0$}
\loigiai{
Để $5x78 < 5178$ thì $x < 1$. Vì $x$ là chữ số nên $x = 0$.
\textbf{Đáp án của câu hỏi là D}
}
\end{ex}

\begin{ex} %[3D1V46]
Trong các số có bốn chữ số khác nhau lập từ bốn chữ số $1; 4; 6; 8$, số lớn thứ hai là
\choice
{\True $8614$}
{$8641$}
{$8461$}
{$8416$}
\loigiai{
Số lớn nhất là 8641. Số lớn thứ hai là 8614.
\textbf{Đáp án của câu hỏi là A}
}
\end{ex}

\begin{ex} %[3D1V46]
Tìm số tự nhiên lớn nhất có bốn chữ số mà tổng các chữ số của nó bằng 12.
\choice
{$9210$}
{$8400$}
{\True $9300$}
{$9120$}
\loigiai{
Để số lớn nhất, hàng nghìn lớn nhất là 9. Còn lại tổng là $12 - 9 = 3$. Đặt 3 ở hàng trăm, hàng chục và đơn vị là 0. Số đó là $9300$.
\textbf{Đáp án của câu hỏi là C}
}
\end{ex}

\textbf{Phần 2. Trắc nghiệm đúng sai}
\setcounterpageref{ex}{0}

\begin{ex} %[3D1V46]
Xét tính đúng sai của việc so sánh các số có bốn chữ số:
\choiceTF
{\True Số có nhiều chữ số hơn thì luôn luôn lớn hơn số có ít chữ số hơn.}
{Hai số có cùng số chữ số thì số nào có chữ số hàng đơn vị lớn hơn sẽ lớn hơn.}
{\True Nếu hai số có hàng nghìn và hàng trăm bằng nhau thì ta so sánh hàng chục.}
{$9876 > 10\,000$.}
\loigiai{
\begin{itemize}
    \item Ý a) Đúng: Quy tắc so sánh số tự nhiên.
    \item Ý b) Sai: Sai vì phải so sánh từ hàng cao nhất trước.
    \item Ý c) Đúng: Thứ tự so sánh chuẩn.
    \item Ý d) Sai: Sai vì 9876 có 4 chữ số, 10 000 có 5 chữ số.
\end{itemize}
\begin{itemize}
    \item a) \textbf{Đúng}.
    \item b) \textbf{Sai}.
    \item c) \textbf{Đúng}.
    \item d) \textbf{Sai}.
\end{itemize}
}
\end{ex}

\begin{ex} %[3D1V46]
Cho bốn số: $4280; 4820; 4082; 4802$:
\choiceTF
{\True Số bé nhất trong bốn số là 4082.}
{\True Số lớn nhất trong bốn số là 4820.}
{Số 4802 lớn hơn số 4820.}
{\True Sắp xếp theo thứ tự giảm dần là: $4820, 4802, 4280, 4082$.}
\loigiai{
\begin{itemize}
    \item Ý a) Đúng: Hàng trăm bé nhất là 0.
    \item Ý b) Đúng: Hàng trăm 8, hàng chục 2.
    \item Ý c) Sai: Sai vì hàng chục $0 < 2$.
    \item Ý d) Đúng: Đúng thứ tự.
\end{itemize}
\begin{itemize}
    \item a) \textbf{Đúng}.
    \item b) \textbf{Đúng}.
    \item c) \textbf{Sai}.
    \item d) \textbf{Đúng}.
\end{itemize}
}
\end{ex}

\textbf{Phần 3. Trả lời ngắn}
\setcounterpageref{ex}{0}

\begin{ex} %[3D1V46]
Tìm số lớn nhất trong các số sau: $6543, 6534, 6453, 6545$.
\shortans{6545}
\loigiai{
\textbf{Đáp số:} $6545$\\
6545 có chữ số hàng đơn vị 5 lớn nhất khi hàng chục, trăm, nghìn bằng nhau.
}
\end{ex}

\begin{ex} %[3D1V46]
Tìm số bé nhất có bốn chữ số khác nhau và có chữ số hàng nghìn bằng 7.
\shortans{7012}
\loigiai{
\textbf{Đáp số:} $7012$\\
Hàng nghìn là 7, tiếp theo lấy nhỏ nhất: 0, 1, 2. Số đó là 7012.
}
\end{ex}

\begin{ex} %[3D1V46]
Có bao nhiêu số tự nhiên ở giữa số 4995 và 5005?
\shortans{9}
\loigiai{
\textbf{Đáp số:} $9$\\
Số lượng số tự nhiên là: $5005 - 4995 - 1 = 9$ số.
}
\end{ex}

\begin{ex} %[3D1V46]
Tìm chữ số $a$ lớn nhất để: $3a52 < 3452$.
\shortans{3}
\loigiai{
\textbf{Đáp số:} $3$\\
Để $3a52 < 3452$ thì $a < 4$. Chữ số lớn nhất là $a = 3$.
}
\end{ex}

\begin{ex} %[3D1V46]
Tìm số tự nhiên chẵn lớn nhất có bốn chữ số khác nhau bé hơn 2000.
\shortans{1986}
\loigiai{
\textbf{Đáp số:} $1986$\\
Hàng nghìn là 1, trăm 9, chục 8, đơn vị chẵn khác nhau lớn nhất là 6. Số đó là 1986.
}
\end{ex}

\begin{ex} %[3D1V46]
Cho dãy số: $3000, 3100, 3200, \dots, 3900$. Dãy số này có bao nhiêu số hạng?
\shortans{10}
\loigiai{
\textbf{Đáp số:} $10$\\
$(3900 - 3000) : 100 + 1 = 10$ số.
}
\end{ex}
